\documentclass[11pt,a4paper]{article}
\usepackage[T1]{fontenc}
\usepackage[utf8]{inputenc}
\usepackage{lmodern,microtype}
\usepackage[a4paper,margin=25mm]{geometry}
\usepackage{amsmath,amssymb,amsthm,mathtools}
\usepackage{booktabs,longtable,array}
\usepackage{enumitem}
\usepackage{hyperref}
\hypersetup{colorlinks=true,linkcolor=blue,citecolor=blue,urlcolor=blue,
 pdftitle={AIM 523: research submission},
 pdfauthor={Siavash Sadeghi; AI-assisted research note}}
\usepackage{fancyhdr}
\pagestyle{fancy}
\fancyhf{}
\fancyhead[L]{\small AIM 523 / Siavash Sadeghi}
\fancyhead[R]{\small Research note / 4 October 2026}
\fancyfoot[C]{\thepage}
\setlength{\headheight}{14pt}
\setlength{\parskip}{4pt}
\setlength{\parindent}{0pt}
\setcounter{tocdepth}{1}
\allowdisplaybreaks
\newtheorem{theorem}{Theorem}[section]
\newtheorem{lemma}[theorem]{Lemma}
\newtheorem{proposition}[theorem]{Proposition}
\theoremstyle{remark}
\newtheorem{remark}[theorem]{Remark}
\newcommand{\R}{\mathbb R}
\newcommand{\N}{\mathbb N}
\newcommand{\eps}{\varepsilon}
\newcommand{\norm}[1]{\left\lVert #1\right\rVert}
\newcommand{\abs}[1]{\left\lvert #1\right\rvert}
\newcommand{\gen}{\mathfrak g}
\newcommand{\dd}{\,\mathrm d}
\newcommand{\sgn}{\operatorname{sgn}}
\newcommand{\supp}{\operatorname{supp}}
\newcommand{\dist}{\operatorname{dist}}
\newcommand{\sech}{\operatorname{sech}}

\begin{document}
\title{A logarithmic lower bound for Nystr\"om estimates}
\author{Siavash Sadeghi}
\date{4 October 2026}
\maketitle
\begin{abstract}
A fixed smooth amplitude on a smooth closed curve with a straight arc gives growth of order square root of log k. This addresses the arbitrary smooth amplitudes and general parametrizations allowed in the target; it does not assert a counterexample for every physical boundary-integral kernel.
\end{abstract}
\textbf{Provenance and review.} Separated and polished from the supplied AI-assisted research note prepared with ChatGPT. Codex reviewed the mathematical argument and accompanying checks. This is a research claim awaiting independent mathematical review; no independent human audit or formal proof verification is claimed.
\medskip
\section{Entry 523: the square-root logarithm is necessary}\label{sec:nystrom}

We use exactly the general-amplitude formulation in catalogue entry 523~\cite{repo523}. The original paper and its fourth version retain the general-curve conjecture~\cite{galkowski}. It is enough to disprove the estimate at $s=0$.

For clarity, the target uses
\[
 L_{1,k}(t,\tau)=k\int_{S^1}e^{ik\langle\gamma(t)-\gamma(\tau),\omega\rangle}
 f_k(\omega,t,\tau)\,\mathrm dS(\omega),\qquad
 \widehat L_{1,k,m}(t)=\frac1{\sqrt{2\pi}}\int_0^{2\pi}e^{-im\tau}L_{1,k}(t,\tau)\,\mathrm d\tau.
\]
All amplitude derivatives must be uniformly bounded in $k$. It suffices to
disprove the asserted bound at $s=0$, where the target quantity is
$k^{-1}\sum_{0<|m|\le(1-\eps)N}\|\widehat L_{1,k,m}\|_{L^2(\Gamma)}$.
The construction uses a single smooth amplitude independent of $k$ and the
admissible choices $\eps=1/2$ and $N=4\lceil k\rceil$.

\subsection{A fixed smooth curve and a fixed amplitude}

Choose a smooth embedded closed curve containing a straight horizontal arc, with a smooth periodic parametrization satisfying, on $|t|<1/4$,
\begin{equation}\label{eq:localcurve}
 \gamma(t)=(\sigma(t),0),\qquad \sigma(t)=t+\frac12t^2.
\end{equation}
Such a curve is obtained by smoothly completing a line segment away from this chart and then making a positive smooth reparametrization. The local speed is $1+t>0$; the global speed is smooth, bounded above, and bounded away from zero. No positive-curvature or unit-speed hypothesis is imposed by the repository target.

Let $\chi$ be a smooth periodic cutoff supported in $(-1/4,1/4)$ and equal to one on $[-1/8,1/8]$. Parametrize a neighbourhood of $(-1,0)\in S^1$ by
\[
 \omega(\theta)=(-\cos\theta,\sin\theta).
\]
Choose a nonnegative even smooth cutoff $\psi$, supported in $|\theta|<\theta_0$ for a sufficiently small fixed $\theta_0>0$, and equal to one near zero. Define, independently of $k$,
\begin{equation}\label{eq:amp}
 f_k(\omega(\theta),t,\tau)=\psi(\theta)\chi(t)\chi(\tau),
\end{equation}
and extend by zero on the rest of the circle. This is a fixed smooth amplitude; every derivative required in the question is uniformly bounded in $k$.

Put $\rho=m/k$ and restrict it to
\[
 J=\left[1-\frac1{32},1+\frac1{32}\right].
\]
For $\theta_0$ small enough, $\tau_\theta=\rho/\cos\theta-1$ lies in the region where $\chi(\tau)=1$, uniformly for $\rho\in J$.

\subsection{Reduction to a nearly degenerate one-dimensional integral}

Using the local coordinate and interchanging the two compact integrals,
\begin{equation}\label{eq:Lhat}
 \widehat L_{1,k,m}(t)=\frac{k\chi(t)}{\sqrt{2\pi}}
 \int\psi(\theta)e^{-ik\sigma(t)\cos\theta}
 \int\chi(\tau)e^{ik[\sigma(\tau)\cos\theta-\rho\tau]}
 \dd\tau\dd\theta.
\end{equation}
The inner phase is exactly quadratic in $\tau$, has critical point $\tau_\theta$, and has second derivative $\cos\theta>0$. Since the cutoff is identically one near the critical point, quadratic stationary phase, with integration by parts away from that point, gives for every $B>0$ a remainder $O(k^{-B})$, uniformly in the parameters. With $u=1+t$, the result is
\begin{equation}\label{eq:stationtau}
 \widehat L_{1,k,m}(t)=\chi(t)\sqrt{k}\,e^{i\pi/4}
 \int \frac{\psi(\theta)}{\sqrt{\cos\theta}}
 e^{ik\Psi(u,\rho,\theta)}\dd\theta+O(k^{-B}),
\end{equation}
where
\begin{equation}\label{eq:Psi}
 \Psi(u,\rho,\theta)=\rho-\frac12u^2\cos\theta
 -\frac{\rho^2}{2\cos\theta}.
\end{equation}
All constants here concern the one fixed curve and amplitude.

Choose a fixed small $\delta_0>0$, for example $1/32$, and take
\[
 t=\rho-1-\delta,\qquad k^{-1/4}\le\delta\le\delta_0.
\]
These $t$ lie in the plateau of $\chi$. Also $u=\rho-\delta<\rho$. The angular phase has just one stationary point on the support, at zero, and
\[
 \Psi_{\theta\theta}(u,\rho,0)=-A,
 \qquad A=\frac{\rho^2-u^2}{2}
       =\frac{\delta(2\rho-\delta)}2\asymp\delta.
\]
The degeneration $A\downarrow0$ as $t\uparrow\rho-1$ is responsible for the logarithm.

\subsection{Uniform stationary phase on the required range}

For completeness, we justify uniformity as $A$ tends to zero with $k$. Let
\[
 g(\theta)=\Psi(u,\rho,0)-\Psi(u,\rho,\theta),\qquad
 b(\theta)=\frac{\psi(\theta)}{\sqrt{\cos\theta}}.
\]
For $|\theta|<\theta_0<\pi/2$,
\begin{equation}\label{eq:gsecond}
 g''(\theta)=\frac{\rho^2}{2\cos\theta}
 -\frac{u^2\cos\theta}{2}
 +\frac{\rho^2\sin^2\theta}{\cos^3\theta}\ge A>0.
\end{equation}
Further,
\[
 g(\theta)=\frac{A\theta^2}{2}+O(\theta^4),
\]
with the remainder uniform in $u,\rho$ in the present compact ranges. Rescale $\zeta=\sqrt{kA}\,\theta$. On each bounded $\zeta$ interval,
\[
 kg\!\left(\frac\zeta{\sqrt{kA}}\right)
 =\frac{\zeta^2}{2}+O\!\left(\frac{\zeta^4}{kA^2}\right),
 \qquad b\!\left(\frac\zeta{\sqrt{kA}}\right)\to1
\]
uniformly, because $kA^2\gtrsim k^{1/2}\to\infty$.

The tails of this rescaled integral are also uniform. On the positive half-line, the derivative of the rescaled phase is increasing and is at least $\zeta$ by~\eqref{eq:gsecond}; the negative half-line is analogous. The amplitudes have uniformly bounded total variation. Integration by parts therefore bounds either tail $|\zeta|>R$ by $C/R$, uniformly in $k,A,u,\rho$. First letting $k\to\infty$ on a bounded interval and then $R\to\infty$ proves
\begin{equation}\label{eq:uniformfresnel}
 \int b(\theta)e^{ik\Psi(u,\rho,\theta)}\dd\theta
 =\frac{e^{ik\Psi(u,\rho,0)}}{\sqrt{kA}}
 \left(\sqrt{2\pi}e^{-i\pi/4}+o(1)\right)
\end{equation}
uniformly for $\rho\in J$ and $k^{-1/4}\le\delta\le\delta_0$.

Combining~\eqref{eq:stationtau} and~\eqref{eq:uniformfresnel} gives
\begin{equation}\label{eq:nystromasymp}
 \widehat L_{1,k,m}(\rho-1-\delta)
 =\sqrt{\frac{2\pi}{A}}\,e^{ik\Psi(u,\rho,0)}(1+o(1)).
\end{equation}
In particular, for all sufficiently large $k$,
\[
 \abs{\widehat L_{1,k,m}(\rho-1-\delta)}^2\ge\frac{c_1}{\delta}
\]
with $c_1>0$ uniform in the specified ranges.

\subsection{The logarithmic lower bound}

On this arc, arclength measure is $\dd s=(1+t)\dd t=u\dd t$, with $u$ bounded above and below by positive constants. Thus, uniformly for integers $m$ with $m/k\in J$,
\begin{align}
 \norm{\widehat L_{1,k,m}}_{L^2(\Gamma)}^2
 &\ge c_2\int_{k^{-1/4}}^{\delta_0}\frac{\dd\delta}{\delta}
 \ge c_3\log k
\end{align}
for sufficiently large $k$. There are at least $c_4k$ such positive integers $m$.

Choose $\eps=1/2$ and $N=4\lceil k\rceil$. The prescribed sum includes all these integers, and at $s=0$ its extra weight is one. Hence
\begin{equation}\label{eq:loglower}
 \frac1k\sum_{0<|m|\le(1-\eps)N}
 \norm{\widehat L_{1,k,m}}_{H_k^0(\Gamma)}
 \ge c_5\sqrt{\log k}\longrightarrow\infty.
\end{equation}

\begin{theorem}
The constant bound in the repository's general smooth-amplitude formulation of entry 523 is false. A fixed smooth curve, a fixed smooth nonconstant-speed parametrization, and a fixed smooth amplitude already force the square-root logarithmic loss.
\end{theorem}

\begin{remark}
The amplitude localization in~\eqref{eq:amp} is explicitly permitted by the repository statement. This argument does not identify that amplitude with every particular physical single- or double-layer kernel. It also does not contradict the positive-curvature, unit-speed result, since the construction has a flat arc and variable speed.
\end{remark}



\begin{thebibliography}{9}
\bibitem{repo523}
M. Colbrook, AIM catalogue, \href{https://github.com/MColbrook/AIM/blob/59a8f0c2957dd272f56bcf282dfe0c066a15f441/problems/523-nystrom-logarithmic-loss.md}{entry 523: Removing the logarithmic loss in Helmholtz Nystr\"om estimates}. Repository statement checked 4 October 2026.

\bibitem{galkowski}
J. Galkowski, M. Rachh, and E. A. Spence, \emph{Helmholtz boundary integral methods and the pollution effect}, \href{https://arxiv.org/abs/2507.22797v4}{arXiv:2507.22797v4} (2026), equations (9.3)--(9.6), Lemmas 11.1 and 11.3, and Conjecture 11.4.
\end{thebibliography}
\end{document}
